% !PRINTED 293
\weberpage{293}

In diesem Falle können wir, da die Vorzeichen von $\alpha$, $\beta$,
$\gamma$, $\delta$ alle gleichzeitig umgekehrt werden dürfen, annehmen:
\[
\alpha+\delta\equiv 2, \quad \alpha\delta-\beta\gamma\equiv 1 \quad (\mathrm{mod}\ p),
\]
und so erhalten wir aus (8):
\[
\begin{aligned}
\alpha_2 &\equiv -1+2\alpha, & \beta_2 &\equiv 2\beta\\
\delta_2 &\equiv -1+2\delta, & \gamma_2 &\equiv 2\gamma\\
\alpha_3 &\equiv -2+3\alpha, & \beta_3 &\equiv 3\beta\\
\delta_3 &\equiv -2+3\delta, & \gamma_3 &\equiv 3\gamma
\end{aligned}\ \ (\mathrm{mod}\ p),
\]
woraus durch den Schluß von $m$ auf $m+1$ gefolgert wird:
\weq{9}{\begin{aligned}
\alpha_m &\equiv -(m-1)+m\alpha, & \beta_m &\equiv m\beta\\
\delta_m &\equiv -(m-1)+m\delta, & \gamma_m &\equiv m\gamma
\end{aligned}\ \ (\mathrm{mod}\ p).}
\gesperrt{Hieraus ersieht man, daß die sämtlichen $A^m$ wieder
zum Falle I. gehören, und daß $p$ der Grad von $A$ ist.}

Die analoge Betrachtung der beiden anderen Fälle ist von
\gesperrt{Serret} durchgeführt, erscheint aber für unseren Zweck entbehrlich.

Dagegen wollen wir hier noch den Satz hinzufügen, \gesperrt{daß
die ganze Gruppe $\mathfrak{L}_0$ sich zusammensetzen läßt aus den
beiden folgenden speziellen Substitutionen}
\weq{10}{A=\begin{pmatrix}1, & 0\\ 1, & 1\end{pmatrix}, \quad B=\begin{pmatrix}0, & 1\\ -1, & 0\end{pmatrix}.}
Der Beweis ergibt sich aus \S\,30, wenn man beachtet, daß
zu jedem der Bedingung $\alpha\delta-\beta\gamma\equiv 1\ (\mathrm{mod}\ p)$ genügenden
Zahlensystem $\alpha$, $\beta$, $\gamma$, $\delta$ sich das Zahlensystem $\alpha'$, $\beta'$, $\gamma'$, $\delta'$ so
wählen läßt, daß
\[
\begin{pmatrix}\alpha', & \beta'\\ \gamma', & \delta'\end{pmatrix} \equiv \begin{pmatrix}\alpha, & \beta\\ \gamma, & \delta\end{pmatrix}\ (\mathrm{mod}\ p) \ \text{und}\ \alpha'\delta'-\beta'\gamma'=1
\]
wird. Es sind also alle Substitutionen von $\mathfrak{L}_0$ unter den in \S\,30
betrachteten enthalten. Wir können hier aber auch leicht die
Zusammensetzung einer beliebigen Substitution in $\mathfrak{L}_0$ aus $A$ und
$B$ wirklich darstellen, und so unabhängig von \S\,30 den Beweis
führen.

Denn wenn zunächst $c$ eine beliebige, durch $p$ nicht teilbare
Zahl ist, so ergibt sich durch wirkliche Ausrechnung leicht
\weq{11}{A^c=\begin{pmatrix}1, & 0\\ c, & 1\end{pmatrix}, \quad C=\begin{pmatrix}c, & 0\\ 0, & c^{-1}\end{pmatrix}=A^{-c^{-1}}BA^{-c}BA^{-c^{-1}}B,}
%CONT>>
